% !TEX program = pdflatex
% =====================================================================
%  GUÍA PARA EL TFG / TFM · recomendaciones de la dirección
%  Compilar en local:  bash tools/build.sh guia/guia_tfg
%  En Overleaf:        Menú > Documento principal > guia/guia_tfg.tex
% =====================================================================
\documentclass[11pt,a4paper]{article}
\newcommand{\idiomadocumento}{spanish}
\newcommand{\estilocitas}{ieee}
% =====================================================================
%  PREÁMBULO COMÚN (memoria, informes y guía)
%  Normalmente NO hace falta tocar este fichero: los datos van en
%  config/datos.tex. Si añades paquetes, hazlo al final, en la zona
%  marcada, y comprueba que el documento sigue compilando sin errores.
%  Requiere: \idiomadocumento y \estilocitas (definidos en datos.tex).
% =====================================================================

% --- Codificación, fuentes y microtipografía --------------------------
\usepackage[T1]{fontenc}
\usepackage[utf8]{inputenc}
\usepackage{lmodern}
\usepackage[babel=true]{microtype}

% --- Idiomas: castellano, catalán e inglés; el principal sale de datos.tex
%     es-tabla        -> «Tabla» en vez de «Cuadro»
%     es-noshorthands -> evita conflictos de < > " con TikZ y tablas
%     es-nodecimaldot -> 3.5 en fórmulas no se convierte en 3,5
\edef\tfgbabelopts{english,catalan,spanish,main=\idiomadocumento}
\expandafter\usepackage\expandafter[\tfgbabelopts,es-tabla,es-noshorthands,es-nodecimaldot]{babel}
\usepackage{csquotes}

% --- Página: A4, márgenes, interlineado 1,15, párrafos sin línea en blanco
\usepackage[a4paper,top=2.5cm,bottom=2.5cm,left=3cm,right=2.5cm,headheight=15pt]{geometry}
\usepackage{setspace}
\setstretch{1.15}
\usepackage[skip=6pt plus 2pt]{parskip}   % separación entre párrafos sin dobles saltos
\raggedbottom
% Sin líneas viudas ni huérfanas
\clubpenalty=10000
\widowpenalty=10000
\displaywidowpenalty=10000

% --- Colores, figuras, tablas y diagramas -----------------------------
\usepackage[table,dvipsnames]{xcolor}
\usepackage{graphicx}
\graphicspath{{figuras/}}
\usepackage{float}
\usepackage[section]{placeins}            % las figuras no "huyen" a otra sección
\usepackage{booktabs,tabularx,multirow,array,longtable}
\newcolumntype{L}{>{\raggedright\arraybackslash}X}
\newcolumntype{C}{>{\centering\arraybackslash}X}
\renewcommand{\arraystretch}{1.2}
\usepackage{amsmath,amssymb}
\usepackage{siunitx}                      % \qty{100}{ms}, \num{12345}: cifras y unidades
\usepackage{tikz}
\usetikzlibrary{arrows.meta,positioning,fit,backgrounds,shapes.geometric,babel}
\usepackage{pgfgantt}
\usepackage{qrcode}
\usepackage{enumitem}
\setlist{itemsep=2pt,topsep=4pt}

% --- Pies de figura y tabla: centrados, breves; tablas con el pie ARRIBA
\usepackage[font=small,labelfont=bf,labelsep=colon,justification=centering,
            singlelinecheck=true]{caption}
\captionsetup[table]{position=top,skip=6pt}
\captionsetup[figure]{skip=6pt}
\usepackage[list=true]{subcaption}

% \fuente{...}: línea de fuente bajo la figura o la tabla (OBLIGATORIA).
%   \fuente{Elaboración propia.}   \fuente{Adaptado de \cite{clave}.}
\newcommand{\fuente}[1]{\par\vspace{3pt}{\centering\footnotesize\fuentename: #1\par}}

% --- Código fuente --------------------------------------------------
\usepackage{listings}
\lstset{
  basicstyle=\ttfamily\small, columns=fullflexible, keepspaces=true,
  breaklines=true, breakatwhitespace=false, tabsize=2,
  frame=tb, framesep=4pt, numbers=left, numberstyle=\tiny\color{gray}, numbersep=6pt,
  keywordstyle=\color{RoyalBlue}\bfseries, commentstyle=\color{ForestGreen}\itshape,
  stringstyle=\color{BrickRed}, showstringspaces=false, captionpos=b,
  floatplacement=htbp,  % con [float] en cada listado, el código no se parte entre páginas
  extendedchars=true,
  literate={á}{{\'a}}1 {é}{{\'e}}1 {í}{{\'i}}1 {ó}{{\'o}}1 {ú}{{\'u}}1
           {Á}{{\'A}}1 {É}{{\'E}}1 {Í}{{\'I}}1 {Ó}{{\'O}}1 {Ú}{{\'U}}1
           {à}{{\`a}}1 {è}{{\`e}}1 {ò}{{\`o}}1 {ï}{{\"i}}1 {ü}{{\"u}}1
           {ñ}{{\~n}}1 {Ñ}{{\~N}}1 {ç}{{\c c}}1 {Ç}{{\c C}}1 {·}{{$\cdot$}}1
           {¿}{{?`}}1 {¡}{{!`}}1
}

% --- Acrónimos: \ac{TFG} -> 1.ª vez «Trabajo de Fin de Grado (TFG)», después «TFG»
\usepackage[printonlyused]{acronym}

% --- Bibliografía (biblatex + biber). Un ÚNICO estilo en todo el documento.
\edef\tfgbibopts{backend=biber,style=\estilocitas}
\expandafter\usepackage\expandafter[\tfgbibopts]{biblatex}
\usepackage{xurl}                          % URLs largas que se parten bien
% Fecha de consulta y URL en el idioma del documento (nunca «Accessed» en un
% texto en castellano):  «Consultado el 15 de nov. de 2026. Disponible en: URL»
\DefineBibliographyStrings{spanish}{urlseen={Consultado el}, url={Disponible en}}
\DefineBibliographyStrings{catalan}{urlseen={Consultat el}, url={Disponible a}}

% --- Textos que dependen del idioma --------------------------------------
\addto\captionsspanish{%
  \def\fuentename{Fuente}\def\acronimosname{Acrónimos}%
  \def\lstlistingname{Listado}\def\lstlistlistingname{Índice de listados}%
  \def\resumenname{Resumen}\def\clavesname{Palabras clave}%
  \sisetup{output-decimal-marker={,}}}
\addto\captionscatalan{%
  \def\fuentename{Font}\def\acronimosname{Acrònims}%
  \def\lstlistingname{Llistat}\def\lstlistlistingname{Índex de llistats}%
  \def\resumenname{Resum}\def\clavesname{Paraules clau}%
  \sisetup{output-decimal-marker={,}}}
\addto\captionsenglish{%
  \def\fuentename{Source}\def\acronimosname{Acronyms}%
  \def\lstlistingname{Listing}\def\lstlistlistingname{List of Listings}%
  \def\resumenname{Abstract}\def\clavesname{Keywords}%
  \sisetup{output-decimal-marker={.}}}

% --- Encabezados de capítulo y profundidad del índice ----------------
\usepackage{titlesec}
\ifdefined\chapter
  \titleformat{\chapter}[display]
    {\normalfont\bfseries\huge\raggedright\hyphenpenalty=10000}
    {\large\chaptertitlename~\thechapter}{6pt}{\huge}
    [\vspace{2pt}\titlerule]
  \titlespacing*{\chapter}{0pt}{0pt}{20pt}
\fi
% Títulos sin palabras partidas por guiones
\titleformat*{\section}{\normalfont\Large\bfseries\raggedright\hyphenpenalty=10000}
\titleformat*{\subsection}{\normalfont\large\bfseries\raggedright\hyphenpenalty=10000}
\setcounter{secnumdepth}{2}   % numera hasta X.Y.Z; no uses un 4.º nivel
\setcounter{tocdepth}{1}      % el índice muestra capítulos y secciones (máx. 2 niveles)

% --- Cabecera y pie: número de página centrado abajo ----------------
\usepackage{fancyhdr}
\pagestyle{fancy}
\fancyhf{}
\fancyhead[R]{\small\nouppercase{\leftmark}}
\fancyfoot[C]{\thepage}
\renewcommand{\headrulewidth}{0.4pt}
\fancypagestyle{plain}{\fancyhf{}\fancyfoot[C]{\thepage}\renewcommand{\headrulewidth}{0pt}}

% --- Enlaces y metadatos del PDF (hyperref casi al final) --------------
\usepackage[colorlinks=true,linkcolor=black,citecolor=black,
            urlcolor=MidnightBlue,bookmarksnumbered=true,pdfusetitle]{hyperref}
\usepackage{bookmark}

% --- Zona para tus paquetes y comandos propios -----------------------
% \usepackage{...}

\usepackage[mode=final]{revnotes}
\addbibresource{guia/guia.bib}
\usepackage[most]{tcolorbox}
\usepackage[titles]{tocloft}
\setlength{\cftsecnumwidth}{2em}
\setlength{\cftsubsecnumwidth}{2.8em}
\setcounter{tocdepth}{2}

\newcommand{\guiaversion}{1.0 (noviembre de 2026)}
\hypersetup{pdftitle={Guía para el TFG/TFM},pdfauthor={Vitor da Silva}}

% --- Cajas -----------------------------------------------------------
\newtcolorbox{regla}[1][]{enhanced, breakable, colback=RoyalBlue!5, colframe=RoyalBlue!70!black,
  fonttitle=\bfseries, title={#1}, boxrule=0.6pt, arc=2pt, left=6pt, right=6pt}
\newtcolorbox{aviso}[1][]{enhanced, breakable, colback=BrickRed!5, colframe=BrickRed!80!black,
  fonttitle=\bfseries, title={#1}, boxrule=0.6pt, arc=2pt, left=6pt, right=6pt}
\newcommand{\cmd}[1]{\texttt{\textbackslash #1}}
\newcommand{\antes}[1]{\textcolor{BrickRed!80!black}{\textbf{Antes:}}~«#1»}
\newcommand{\despues}[1]{\textcolor{ForestGreen!80!black}{\textbf{Después:}}~«#1»}
\newlist{checklist}{itemize}{1}
\setlist[checklist]{label=$\square$, itemsep=1pt, leftmargin=1.6em}

\begin{document}

% =====================================================================
\begin{titlepage}
  \centering
  \vspace*{3cm}
  {\Huge\bfseries Guía para el TFG y el TFM\par}
  \vspace{0.8cm}
  {\Large Estructura, redacción, citas, figuras, pruebas y uso de la IA\par}
  \vspace{2.5cm}
  {\large Grau en Enginyeria Informàtica\\ Escola Politècnica Superior · Universitat de Lleida\par}
  \vspace{2.5cm}
  {\large Vitor da Silva\par}
  \vfill
  {\small Versión \guiaversion\par}
\end{titlepage}

\tableofcontents
\clearpage

% =====================================================================
\section*{Antes de empezar}
\addcontentsline{toc}{section}{Antes de empezar}

Esta guía reúne las recomendaciones que doy a todos los estudiantes que hacen el
Trabajo de Fin de Grado (TFG) o de Máster (TFM) conmigo. Sale de la normativa de
la EPS~\cite{eps_reglament2024,eps_portafoli2025,eps_terminis2025} y, sobre todo,
de los comentarios que he repetido una y otra vez al revisar memorias. Casi
todos los errores que verás aquí los cometió alguien antes que tú; evitarlos
desde el principio te ahorra semanas de correcciones al final.

Úsala junto con la plantilla de LaTeX: la plantilla ya aplica la mayoría de las
reglas de formato, y cada capítulo trae comentarios (\texttt{\% GUÍA:}) con lo
que tiene que contener. Lee la guía entera una vez al empezar y vuelve a las
listas de comprobación de la Sección~\ref{sec:checklists} antes de cada entrega.

\begin{regla}[Las diez reglas que más nota cuestan]
\begin{enumerate}[leftmargin=1.5em]
  \item \textbf{Trabajo propio.} Nunca uses el TFG de otra persona como base: se
        coteja y un solapamiento grande hace que el trabajo no sea defendible.
  \item \textbf{Portada oficial completa}: título idéntico al de la propuesta,
        todos los directores, mes y año. Es el fallo más repetido.
  \item \textbf{Objetivos numerados y medibles} (O1, O2\ldots), que se cierran
        uno a uno en las conclusiones.
  \item \textbf{Trazabilidad}: objetivo $\rightarrow$ requisito $\rightarrow$
        prueba $\rightarrow$ conclusión, con una tabla que lo demuestre.
  \item \textbf{Lo que no has hecho no se presenta como hecho.} Una prueba
        propuesta no es una prueba realizada.
  \item \textbf{Pruebas con evidencia}: casos con resultado esperado y obtenido,
        no «se ha probado manualmente».
  \item \textbf{Toda afirmación con fuente} y un único estilo de cita en todo el
        documento.
  \item \textbf{Resumen, Resum y Abstract}, cada uno en su página y con el mismo
        contenido.
  \item \textbf{Un único cronograma}, en pasado, que llegue hasta el depósito y
        la defensa.
  \item \textbf{Toda figura y tabla citada e interpretada en el texto}, con pie
        breve y fuente.
\end{enumerate}
\end{regla}

% =====================================================================
\section{El proceso: calendario y evaluación}
\label{sec:proceso}

\subsection{Fases y plazos}

El TFG del grado es un trabajo individual y autónomo de 15~ECTS, unas 375 horas
de dedicación, hecho bajo la orientación de la dirección. La
Tabla~\ref{tab:fases} resume las fases según el reglamento de la
EPS~\cite{eps_reglament2024} y el esquema de plazos~\cite{eps_terminis2025}. Las
fechas concretas cambian cada curso: consúltalas siempre en la web de la EPS.

\begin{table}[htbp]
  \centering
  \caption{Fases del TFG/TFM y qué hay que entregar en cada una.}
  \label{tab:fases}
  \small
  \begin{tabularx}{\textwidth}{@{}p{3.1cm}p{3.6cm}L@{}}
    \toprule
    \textbf{Fase} & \textbf{Cuándo} & \textbf{Qué se entrega o se exige} \\
    \midrule
    Propuesta & Al menos 10 días antes del plazo de matrícula & Documento de propuesta con director/es, título previsto, objetivos, temporización y palabras clave. Haber superado el 75\,\% de los créditos (grado de 4 años) o el 40\,\% (máster). \\
    Matrícula & Inicio del 1.\textsuperscript{er} o del 2.\textsuperscript{o} semestre & Propuesta ya presentada y matrícula de todos los créditos pendientes. \\
    Informe inicial & Semanas 4--6 & Objetivos, índice previsto, tareas, planificación y fuentes consultadas. 10\,\% de la nota. \\
    Informe de seguimiento & Semanas 10--12 & Evolución frente al plan, evaluación crítica, decisiones técnicas y planificación de lo pendiente. 10\,\% de la nota. \\
    Documento final & Al menos 3 semanas antes de la defensa & La memoria completa. 50\,\% de la nota. \\
    Depósito & Al menos 5 días hábiles antes de la defensa & Solo si los tres ítems anteriores suman 3,5 puntos o más. \\
    Defensa & En las fechas que publica la EPS & Presentación pública de 45 minutos como máximo y preguntas del tribunal. 30\,\% de la nota. \\
    \bottomrule
  \end{tabularx}
  \fuente{Elaboración propia a partir de~\cite{eps_reglament2024,eps_terminis2025}.}
\end{table}

Las plantillas de \texttt{informes/} cubren exactamente lo que el Portafoli
pide para el informe inicial y el de seguimiento. Entrégalos a tiempo: son el
20\,\% de la nota y la única ocasión de corregir el rumbo antes del final.

\subsection{Qué se evalúa en la memoria}

El documento final pesa la mitad de la nota y lo evalúa la dirección con el
Portafoli~\cite{eps_portafoli2025}. Conviene conocer los criterios antes de
escribir:

\begin{itemize}
  \item \textbf{Redacción y estructura (0--2 puntos):} memoria entendible y
        coherente, ortografía y gramática correctas, contenido original y uso
        de herramientas adecuadas para documentos científico-técnicos.
  \item \textbf{Dificultad e innovación (0--1):} innovación, integración de
        conocimientos del grado, aprendizaje de tecnologías nuevas y pruebas de
        laboratorio o de campo.
  \item \textbf{Desarrollo (0--2):} análisis y diseño, planificación, calidad
        del resultado, \textbf{estudio de viabilidad económica y legal},
        documentos técnicos (código, resultados), uso de fuentes especializadas,
        capacidad de análisis y síntesis, y \textbf{contexto socioeconómico y
        sostenibilidad}.
\end{itemize}

Los aspectos de forma (portada, resúmenes, figuras, citas) son rápidos de
corregir y suben nota; no los dejes para la última semana.

\subsection{El depósito}

Para depositar se envía, desde el correo de la UdL y con copia a la dirección,
a la secretaría del campus (Igualada: \texttt{igualada.academic@udl.cat}; Lleida:
\texttt{eps.tfg-tfm@udl.cat})~\cite{eps_terminis2025}:

\begin{itemize}
  \item la memoria en PDF con la portada oficial;
  \item el resumen del TFG/TFM en PDF, con la plantilla oficial;
  \item los anexos o el código, comprimidos o como enlace a la nube;
  \item firmados digitalmente: el documento 2-1 (Biblioteca), el 2-2
        (autorización para el repositorio RECERCAT) y el 2-3 (premios).
\end{itemize}

El nombre y los apellidos del estudiante van en el asunto del correo y en el
nombre de cada fichero. Los envíos cuentan de lunes a las 9:00 a viernes a las
14:00. En los grados, además, hay que tener acreditada la tercera lengua.

% =====================================================================
\section{Cómo trabajar con la plantilla}
\label{sec:plantilla}

\begin{enumerate}
  \item Crea tu repositorio \textbf{privado} a partir de la plantilla en GitHub
        (botón \emph{Use this template}) e impórtalo en Overleaf (\emph{New
        Project $\rightarrow$ Import from GitHub}). Comparte el proyecto con la
        dirección desde el primer día.
  \item Rellena \texttt{config/datos.tex}: autor, título, dirección, idioma
        (castellano, catalán o inglés) y estilo de cita. Todo lo demás sale de
        ahí.
  \item Descarga la portada oficial de la EPS, rellénala, expórtala a PDF y
        guárdala como \texttt{portada/portada.pdf}.
  \item Escribe un capítulo por fichero en \texttt{capitulos/}. Borra el texto
        de ejemplo a medida que escribes el tuyo.
  \item Para lo que tengas pendiente, usa \cmd{pending}\{\ldots\}: sale en rojo
        mientras trabajas. La dirección te dejará notas con
        \cmd{vitornote}\{\ldots\}.
  \item Antes de entregar, cambia \texttt{mode=review} por \texttt{mode=final}
        en \texttt{main.tex}. Desaparecen las notas, y si queda algún
        \cmd{pending} o falta la portada oficial, la compilación falla: así no
        se deposita nada provisional.
\end{enumerate}

Algunas herramientas de la plantilla que conviene usar desde el principio:
\cmd{ac}\{API\} para las siglas (la primera vez se escribe completa),
\cmd{fuente}\{\ldots\} bajo cada figura y tabla, \texttt{Figura\textasciitilde\cmd{ref}\{fig:x\}}
para citar figuras y \cmd{cite}\{clave\} para las referencias. Nunca escribas a
mano el número de una figura, una tabla o una referencia.

% =====================================================================
\section{Estructura de la memoria}
\label{sec:estructura}

\subsection{Parte preliminar}

El orden es: portada oficial; Resumen (castellano), Resum (catalán) y Abstract
(inglés), cada uno en su página, con el mismo contenido y con palabras clave;
agradecimientos (opcional); declaración de uso de IA; índices de figuras, de
tablas y de listados; lista de acrónimos, y por último el índice general, con
dos niveles como máximo. Si una página se titula «Abstract», el texto está en
inglés.

El resumen no es una introducción: dice qué problema hay, qué has hecho, qué
has obtenido (con cifras) y qué aporta. Quien lo lea debe conocer el resultado.

\subsection{Capítulos}

La Tabla~\ref{tab:capitulos} resume la función de cada capítulo de la plantilla.
La estructura se puede adaptar al tipo de trabajo, pero cada capítulo tiene que
tener una función clara, empezar con una frase que la explique y estar conectado
con los objetivos. Si la introducción anuncia un capítulo, ese capítulo existe.

\begin{table}[htbp]
  \centering
  \caption{Función de cada capítulo de la memoria.}
  \label{tab:capitulos}
  \small
  \begin{tabularx}{\textwidth}{@{}p{3.6cm}L@{}}
    \toprule
    \textbf{Capítulo} & \textbf{Qué debe contener} \\
    \midrule
    1. Introducción & Contexto, motivación (sin adelantar resultados), objetivos O1\ldots On medibles, alcance y estructura del documento. \\
    2. Fundamentos y estado del arte & Los conceptos que necesita el lector, explicados y con un elemento visual, y una comparativa con criterios de las soluciones existentes (estudio de mercado en una app). Termina con el hueco que cubre tu trabajo. \\
    3. Metodología, planificación y viabilidad & Metodología justificada, un único cronograma hasta la defensa, desviaciones, herramientas con versión, viabilidad económica y legal, y sostenibilidad. \\
    4. Análisis y diseño & Requisitos RF/RNF con ID y verificables, arquitectura explicada y decisiones justificadas; en una app, prototipos y mapa de navegación. \\
    5. Implementación & Estructura del código, componentes, fragmentos clave explicados, capturas reales, seguridad y gestión de errores. \\
    6. Pruebas y resultados & Casos de prueba, resultados interpretados, tabla de trazabilidad, discusión y limitaciones. \\
    7. Conclusiones y trabajo futuro & Cumplimiento objetivo a objetivo, conclusiones, limitaciones y $\sim$1 página de trabajo futuro. \\
    Apéndices & Repositorio con la versión entregada, vídeo demo con QR (apps) y manual de instalación. \\
    \bottomrule
  \end{tabularx}
  \fuente{Elaboración propia.}
\end{table}

\subsection{Según el tipo de trabajo}

\paragraph{Aplicaciones móviles o web.} Enseña la aplicación: estudio de mercado
con las mismas fichas para todas las alternativas (con logo o captura y la fecha
de consulta de los datos de la tienda), prototipos, mapa de navegación y
capturas reales explicadas. Graba un \textbf{vídeo demo} en un dispositivo real,
súbelo a YouTube (puede ser «no listado») e inclúyelo en la memoria con su
enlace y un código QR. La memoria tiene que entenderse igualmente sin el vídeo.

\paragraph{Sistemas, cloud y datos o aprendizaje automático.} Describe el
entorno experimental de forma reproducible (hardware, versiones, configuración).
Para cada dataset: estructura, número de muestras y particiones, tipo de dato
con un ejemplo visual, referencia y protocolo de evaluación. Define las métricas
con su fórmula, explica cada hiperparámetro y da los resultados en el punto de
operación real, no el mejor caso.

\paragraph{TFM.} Todo lo anterior, y además: una sección que diga sin rodeos qué
aporta el trabajo frente al estado del arte, preguntas u objetivos de
investigación, un estado del arte crítico con fuentes recientes, comparación con
una referencia (\emph{baseline}) y rigor estadístico cuando se afirme que una
diferencia es significativa.

% =====================================================================
\section{Redacción y estilo}
\label{sec:redaccion}

\begin{itemize}
  \item \textbf{Una idea por frase.} Por encima de unas 40 palabras, divide.
  \item \textbf{Prosa, no viñetas.} Las listas son para enumeraciones cortas y
        siempre llevan una frase que las introduce. No hay listas de un solo
        elemento.
  \item \textbf{Todo apartado numerado lleva texto} antes de su primer
        subapartado. Si no tienes nada que decir entre 4.5 y 4.5.1, no numeres.
        No uses un cuarto nivel de títulos (4.5.1.1).
  \item \textbf{Tiempos verbales coherentes.} Lo que ya has hecho, en pasado:
        \antes{El desarrollo se llevará a cabo entre septiembre y febrero} $\rightarrow$
        \despues{El desarrollo se llevó a cabo entre septiembre y julio}. El
        futuro solo en el trabajo futuro (y el «deberá» de los requisitos).
  \item \textbf{Negrita solo para etiquetas} de una lista o definición seguidas
        de dos puntos, nunca para resaltar palabras en mitad de una frase. Para
        un énfasis puntual, cursiva.
  \item \textbf{Siglas y nombres propios definidos en su primera aparición}:
        «Trabajo de Fin de Grado (TFG)». El nombre de tu aplicación se presenta
        en la introducción, no en la conclusión.
  \item \textbf{Explica los términos técnicos} la primera vez (\emph{batch size},
        \emph{token}, latencia\ldots): tú los conoces, el tribunal quizá no.
  \item \textbf{Sin absolutos ni sobreafirmaciones.} \antes{la única solución
        viable} $\rightarrow$ \despues{una de las soluciones más adecuadas para
        este caso}. Las hipótesis no son hechos: \antes{se confirma el
        sobreajuste} $\rightarrow$ \despues{hay indicios de sobreajuste}.
  \item \textbf{Un solo idioma.} Los rótulos de figuras y tablas, la interfaz que
        describes y la bibliografía siguen el idioma del documento («Tabla 1», no
        «Table 1»). Los términos en otro idioma, en cursiva.
  \item \textbf{Pasa un corrector} (LanguageTool o el de tu editor) antes de cada
        entrega, también sobre el texto de las imágenes.
\end{itemize}

% =====================================================================
\section{Citas y bibliografía}
\label{sec:citas}

\subsection{Qué necesita una cita}

Toda afirmación que no sea conocimiento común o resultado tuyo: datos
estadísticos y de mercado, cifras de usuarios o descargas, afirmaciones de
rendimiento («más rápido», «más escalable»), características de una tecnología,
normativa, definiciones y comparativas. La cita va pegada a la afirmación que
sostiene, no al final del párrafo.

\begin{center}
  \antes{Diversos estudios han demostrado que\ldots} $\rightarrow$
  \despues{García y Ruiz~[2] mostraron que\ldots}
\end{center}

\subsection{Estilo}

Usa \textbf{un único estilo} en todo el documento; recomiendo el numérico de
IEEE (\texttt{[1]}), que es el que trae la plantilla. Si prefieres autor-año
(APA), cámbialo en \texttt{config/datos.tex}, pero no mezcles: no vale citar
con números y tener la bibliografía en formato autor-año, ni citar por el nombre
de la revista en lugar del autor. Con biblatex, el formato de cada entrada sale
solo; tu trabajo es que los datos de la entrada sean correctos.

\subsection{Calidad de las fuentes}

\begin{itemize}
  \item Prefiere fuentes primarias u oficiales: artículos, libros, normas, RFC,
        \emph{white papers} y la documentación oficial de cada herramienta.
  \item Wikipedia, blogs, vídeos, foros o Pinterest no sirven como fuente
        principal. Si citas un \emph{preprint} de arXiv que ya está publicado,
        cita la versión publicada.
  \item \textbf{Verifica cada entrada contra el original}: autores, título,
        año, revista, DOI. Una referencia con la autoría equivocada o un
        artículo «enviado» citado como publicado es un error grave.
  \item La fecha de consulta solo se pone en las páginas web, y en el idioma
        del documento («Consultado el\ldots», no «Accessed» en un texto en
        castellano). Los datos que cambian (valoraciones de tienda, número de
        descargas) siempre con su fecha.
  \item Cada entrada de la bibliografía se cita en el texto y cada cita está en
        la bibliografía. Con biblatex esto es automático si nunca escribes la
        bibliografía a mano.
  \item Como orientación, un TFG sólido suele tener entre 20 y 30 referencias de
        calidad; con 8 o 9 el estado del arte se queda corto.
  \item Cada librería, herramienta, método o dataset que uses lleva su
        referencia y su versión exacta.
\end{itemize}

La biblioguía de la UdL~\cite{udl_biblioguia} explica cómo buscar fuentes y
gestionar referencias con Zotero, que puede exportar directamente a
\texttt{bibliografia.bib}.

% =====================================================================
\section{Figuras, tablas y código}
\label{sec:figuras}

\begin{regla}[Regla de oro de los elementos visuales]
Todo elemento visual se cita en el texto, a ser posible antes de que aparezca, y
el texto guía al lector por él: qué mirar, qué significa y qué se concluye. «Los
resultados están en la Tabla 5» no basta.
\end{regla}

\begin{center}
  \antes{Los resultados de las pruebas están en la Tabla 5.}\\[2pt]
  $\downarrow$\\[2pt]
  \despues{La Tabla 5 recoge las 12 pruebas de integración: 11 se superaron; la
  de reenvío con pérdida de paquetes falló, lo que motivó el reintento descrito
  en la Sección 7.2.}
\end{center}

\begin{itemize}
  \item \textbf{Numeración automática} con \cmd{label} y \cmd{ref}: continua,
        sin huecos y en orden de aparición. Una tabla es una tabla, no una
        figura.
  \item \textbf{Pie breve y descriptivo} (no «Figura 1» ni el nombre de la
        pantalla sin más), centrado y sin doble puntuación. La explicación
        larga va en el texto. Las tablas llevan el pie arriba; las figuras,
        abajo.
  \item \textbf{Fuente siempre}, con \cmd{fuente}: «Elaboración propia» o
        «Adaptado de [n]». Una imagen ajena sin fuente es un problema de
        integridad.
  \item \textbf{Gráficas a ancho completo}, en formato vectorial (PDF) generado
        con un script que guardas en el repositorio, con ejes etiquetados y
        unidades.
  \item \textbf{Capturas legibles}, agrupadas de dos en dos o de tres en tres
        con \texttt{subfigure}, nunca estiradas ni borrosas. Ninguna figura
        tapa el número de página.
  \item \textbf{Revisa la ortografía dentro de las imágenes}: un «FEBERERO» en
        un diagrama se ve igual que en el texto.
  \item \textbf{Diagramas propios} de la arquitectura y del flujo de datos. Un
        collage de logos de tecnologías no explica nada: conviértelo en una
        tabla de herramientas con versión y motivo.
  \item \textbf{Apoyo visual en la teoría}: cada concepto denso del estado del
        arte gana con un esquema (por ejemplo, las capas de una red neuronal o
        el intercambio de mensajes de un protocolo).
  \item \textbf{Código solo cuando aporta}: fragmentos de 10 a 25 líneas,
        explicados en el texto (qué problema resuelven, qué líneas importan). El
        código completo va en el repositorio.
\end{itemize}

% =====================================================================
\section{Pruebas, resultados y trazabilidad}
\label{sec:pruebas}

\begin{itemize}
  \item \textbf{Evidencia, no intención.} Cada requisito tiene al menos un caso
        de prueba con pasos o entrada, resultado esperado, resultado obtenido y
        estado. Si dices que hay una cobertura del 90\,\% o 150 casos, enséñalo.
  \item \textbf{Lo propuesto no es lo realizado.} Si una prueba con usuarios no
        se ha podido hacer, se dice así y se lleva a trabajo futuro. Cada
        funcionalidad tiene un único estado en todo el documento (hecha,
        parcial o pendiente).
  \item \textbf{Cifras coherentes} entre el resumen, el texto, las tablas y las
        conclusiones. Revísalas todas antes de cada entrega.
  \item \textbf{Pruebas con personas}: indica el número de participantes, su
        perfil, el instrumento (por ejemplo, el cuestionario SUS) y los valores
        absolutos junto a los porcentajes: \antes{el 54,5\,\% prefiere\ldots}
        $\rightarrow$ \despues{el 54,5\,\% (6 de 11) prefiere\ldots}. Con datos
        personales, y sobre todo de menores, explica el consentimiento, la
        anonimización y el cumplimiento del RGPD.
  \item \textbf{Trazabilidad}: una tabla que relacione cada objetivo con sus
        requisitos, sus pruebas y el resultado. Es lo primero que mira el
        tribunal para saber si has hecho lo que prometías.
  \item \textbf{Conclusiones objetivo a objetivo}: cumplido, parcialmente
        cumplido o no cumplido, con la evidencia. Las limitaciones van aparte
        del trabajo futuro, que ocupa alrededor de una página.
  \item \textbf{El título no promete más de lo entregado.} Si el alcance cambia,
        se ajustan el título (antes del depósito), los objetivos y la
        introducción.
\end{itemize}

% =====================================================================
\section{Uso de la inteligencia artificial}
\label{sec:ia}

Las herramientas de IA generativa pueden ayudarte en tareas auxiliares, pero no
sustituyen tu aportación personal, tu reflexión crítica ni tu criterio técnico.
El TFG evalúa lo que sabes hacer tú. La Tabla~\ref{tab:ia} resume qué usos se
admiten.

\begin{table}[htbp]
  \centering
  \caption{Usos de la IA generativa en el TFG/TFM.}
  \label{tab:ia}
  \small
  \begin{tabularx}{\textwidth}{@{}p{4.6cm}p{3.2cm}L@{}}
    \toprule
    \textbf{Uso} & \textbf{¿Se admite?} & \textbf{Condiciones} \\
    \midrule
    Corrección lingüística y de formato & Sí, con declaración & Como apoyo limitado; el texto y las ideas son tuyos. \\
    Generación o revisión de código & Sí, con declaración & Debes entender, probar y poder defender cada línea. \\
    Búsqueda inicial de fuentes & Sí & Cada fuente se lee y se verifica en el original antes de citarla. \\
    Diagramas y visualizaciones & Sí, con declaración & Indica la herramienta en la fuente de la figura. \\
    Análisis, discusión y conclusiones & No & Son la parte que demuestra tu criterio: tienen que ser tuyas. \\
    Redactar capítulos completos & No & Es contenido no original y afecta a la nota de redacción. \\
    \bottomrule
  \end{tabularx}
  \fuente{Elaboración propia.}
\end{table}

\begin{aviso}[Riesgos que veo con frecuencia]
\begin{itemize}[leftmargin=1.2em]
  \item \textbf{Referencias inventadas.} Los asistentes generan citas con
        autores, títulos y DOI que no existen. Cada referencia se comprueba en
        la fuente original; una referencia falsa es un problema de integridad.
  \item \textbf{Texto genérico sin hilo}: viñetas por todas partes, apartados
        numerados sin texto, frases de relleno («en el panorama actual\ldots»)
        y afirmaciones sin fuente. Se nota, y baja la nota de redacción.
  \item \textbf{Reestructuraciones masivas} propuestas por un asistente cuando el
        documento ya está casi cerrado: consulta antes con la dirección, porque
        pueden desestabilizar una memoria que funcionaba.
  \item \textbf{Datos que no son tuyos.} No pegues en herramientas externas datos
        personales, datos de terceros ni código o documentación confidencial de
        una empresa.
\end{itemize}
\end{aviso}

Todo uso se declara en la página «Declaración de uso de herramientas de IA» de
la plantilla: qué herramienta (y versión), para qué tarea concreta, en qué partes
y cómo revisaste el resultado. Si no has usado ninguna, también se dice. Tras
usarlas, eres responsable de todo el contenido, incluidos sus errores.

% =====================================================================
\section{Integridad académica}
\label{sec:integridad}

\begin{itemize}
  \item \textbf{No uses el TFG de otra persona como plantilla}, ni su portada,
        ni su estructura de resultados. Las memorias se cotejan entre sí. En un
        caso reciente, una memoria con dos terceras partes de frases idénticas
        a otro TFG se consideró no defendible hasta que se reescribió por
        completo. Para empezar, usa esta plantilla.
  \item \textbf{Parafrasear no exime de citar.} Si la idea es de otra persona, se
        cita aunque la escribas con tus palabras. Las citas literales van entre
        comillas y con referencia.
  \item \textbf{Deja claro qué has hecho tú.} Si partes de un proyecto previo, de
        código de un grupo de investigación o de una empresa, distingue en el
        texto tu aportación de la de los demás.
  \item \textbf{El código ajeno} se usa respetando su licencia y se cita.
\end{itemize}

% =====================================================================
\section{Listas de comprobación}
\label{sec:checklists}

\subsection{Antes del informe inicial}
\begin{checklist}
  \item Objetivos numerados y medibles.
  \item Índice previsto con los capítulos de la plantilla adaptados a tu trabajo.
  \item Tareas con su entregable y cronograma hasta el depósito y la defensa.
  \item Fuentes que de verdad has leído, citadas con \cmd{cite}.
\end{checklist}

\subsection{Antes del informe de seguimiento}
\begin{checklist}
  \item Tabla de control: qué está hecho, qué está en curso y qué no ha empezado.
  \item Desviaciones explicadas y cronograma actualizado.
  \item Decisiones técnicas con las alternativas descartadas y el motivo.
  \item Valoración crítica honesta del estado y de los riesgos.
\end{checklist}

\subsection{Antes del documento final y del depósito}
\begin{checklist}
  \item Portada oficial completa: título idéntico al de la propuesta, todos los
        directores, mes (en catalán) y año.
  \item Resumen, Resum y Abstract, cada uno en su página, con resultados y
        palabras clave.
  \item Declaración de uso de IA rellenada.
  \item La estructura que anuncia la introducción coincide con el índice.
  \item Cada objetivo se cierra en las conclusiones con su evidencia.
  \item Tabla de trazabilidad completa; nada propuesto aparece como realizado.
  \item Toda figura y tabla citada e interpretada, con pie y fuente; sin
        huecos de numeración.
  \item Un único estilo de cita; cada entrada verificada contra el original.
  \item Cronograma en pasado que llega al depósito y la defensa.
  \item Viabilidad económica y legal, y sostenibilidad.
  \item Repositorio con la versión entregada (tag o commit); en una app, vídeo
        demo con enlace y QR que funcionan.
  \item Sin texto provisional: compila sin errores en \texttt{mode=final}, sin
        referencias ni citas indefinidas (\texttt{??}).
  \item Corrector pasado, también sobre el texto de las imágenes.
  \item Ficheros con tu nombre y apellidos, y el resto de documentos del
        depósito firmados.
\end{checklist}

% =====================================================================
\section{Errores frecuentes, con ejemplos}
\label{sec:errores}

La Tabla~\ref{tab:errores} recoge correcciones reales, anonimizadas, de memorias
que he revisado. Si te reconoces en alguna, corrígela antes de entregar.

{\small
\begin{longtable}{@{}p{3.0cm}p{5.8cm}p{5.8cm}@{}}
  \caption{Errores frecuentes y cómo corregirlos.}\label{tab:errores}\\
  \toprule
  \textbf{Tema} & \textbf{Antes} & \textbf{Después} \\
  \midrule
  \endfirsthead
  \toprule
  \textbf{Tema} & \textbf{Antes} & \textbf{Después} \\
  \midrule
  \endhead
  \bottomrule
  \endfoot
  Portada & «Mes: Elija un elemento.» o «Mes: July» & «Mes: Juliol», y todos los directores \\
  Resúmenes & Páginas «Abstracte» y «Abstracto»; una página «Abstract» con el texto en castellano & «Resum», «Resumen» y «Abstract», cada una en su idioma \\
  Objetivos & «Conseguir un rendimiento óptimo» & «O3. Latencia media $<$ 100 ms con 200 usuarios» \\
  Estructura & La introducción anuncia un capítulo de pruebas que no existe & Crear el capítulo o quitarlo del anuncio \\
  Cronograma & Plan de septiembre a febrero; el texto habla de 4 etapas y el Gantt tiene 6 & Un único plan de septiembre a julio, con depósito y defensa \\
  Requisitos & «RNF-02: El sistema será escalable» & «RNF-02: soporta N usuarios con latencia $<$ X ms», o marcarlo como objetivo de diseño no medido \\
  Normativa & «La aplicación cumple WCAG 2.1 AA» (sin medirlo) & «Se han seguido las pautas de WCAG 2.1 [n]», o aportar las medidas \\
  Justificación & «Se eligió porque la usan grandes empresas» & «Se eligió porque resuelve el descubrimiento de nodos y el NAT» \\
  Stack & Las tecnologías se describen en tres capítulos distintos & Una tabla «capa -- herramienta -- versión -- motivo» \\
  Métricas & La métrica principal se explica solo con palabras & La fórmula, p.\,ej. $\mathrm{CER} = (S + D + I)/N$, y qué es cada término \\
  Citas & «(Psicología: Teoría e Pesquisa, 2024)» & «Pereira [4]» o «(Pereira, 2024)», según el estilo elegido \\
  Fuentes & Wikipedia, un vídeo o un foro como fuente de una definición & La documentación oficial, el RFC o el artículo original \\
  Bibliografía & 52 de 66 entradas nunca se citan & Citar cada entrada o eliminarla \\
  Fecha de consulta & «[Accessed 11 Oct. 2025]» en un texto en castellano & «Consultado el 11 de oct. de 2025» \\
  Figuras & Pie sin mención en el texto & «Como muestra la Figura 3, \ldots» y dos o tres frases de lectura \\
  Referencias & «en la 33» o «Error! Reference source not found» & «en la Figura 33», con referencias automáticas \\
  Tabla o figura & «Figura 15» para una tabla; «Table 1» en un texto en castellano & «Tabla 3», con su índice de tablas \\
  Pies & «Figura 7: :», «Table 6.20.:» & «Figura 7:», «Table 6.20:» \\
  Imágenes & «FEBERERO», «JAUSTES FINALES» dentro de un diagrama & «FEBRERO», «AJUSTES FINALES» \\
  Negrita & «una \textbf{motivación social y ética}, basada en\ldots» & «una motivación social y ética, basada en\ldots» \\
  Siglas & «TFG» sin definir; el nombre de la app aparece por primera vez en las conclusiones & «Trabajo de Fin de Grado (TFG)»; la app se presenta en la introducción \\
  Gramática & «que la imagen a cambiado», «se ha realizado las revisiones» & «que la imagen ha cambiado», «se han realizado las revisiones» \\
  Pruebas & «Cobertura total, más de 150 casos» & El porcentaje real de cobertura y una tabla de casos con su estado \\
  Validación & «Se realizaron pruebas con usuarias» (no se hicieron) & «La validación con usuarias no se ha podido realizar; se propone como trabajo futuro» \\
  Resultados & «Una precisión cercana al 90\,\%» & «79,84\,\% (FAR 1\,\%, FRR 42\,\%) en el umbral desplegado» \\
  Precisión técnica & «El JWT se genera mediante bcrypt» & «Las contraseñas se guardan con bcrypt; el JWT se firma con HS256» \\
  Encuestas & «El 54,5\,\% prefiere\ldots» & «El 54,5\,\% (6 de 11) prefiere\ldots», con N, perfil e instrumento \\
  Texto provisional & «QUEDA EL MAPA», «final URL is still to be defined» & El contenido definitivo, o \cmd{pending} mientras se trabaja \\
  Trabajo futuro & Cuatro páginas, o dentro de un anexo & Alrededor de una página al final de las conclusiones \\
\end{longtable}}

% =====================================================================
\section*{Recursos}
\addcontentsline{toc}{section}{Recursos}

\begin{itemize}
  \item Web de la EPS: normativa, plazos, portada oficial y documentos del
        depósito (\url{https://www.eps.udl.cat/}).
  \item Biblioguía de TFG/TFM de la UdL~\cite{udl_biblioguia}.
  \item Documentación de Overleaf: \url{https://www.overleaf.com/learn}.
  \item Gestor de referencias Zotero: \url{https://www.zotero.org/}.
  \item Corrector LanguageTool: \url{https://languagetool.org/}.
\end{itemize}

\printbibliography[heading=bibintoc,title={Referencias}]

\end{document}
